\chapter{Strapdown Propagation of the Nominal State}
\label{ch:strapdown}

\begin{objectives}
\begin{itemize}
  \item State the continuous-time nominal kinematics and understand each term.
  \item Derive the exact solution of the kinematics for constant inputs over a step.
  \item Compare Euler, midpoint and Runge--Kutta integration, and predict their convergence orders.
  \item Propagate to arbitrary timestamps, handle gaps, and initialize the nominal state.
  \item Understand IMU preintegration as used by optimization-based estimators.
\end{itemize}
\end{objectives}

A \emph{strapdown} inertial system has its sensors rigidly attached (strapped down) to the vehicle, as
opposed to mounted on a stabilized gimbal platform. All the work of keeping track of orientation is done in
software, by integrating the gyroscope. This chapter is about that integration.

% ------------------------------------------------------------------------------
\section{Continuous-time kinematics}

Taking the IMU model \eqref{eq:imu-model}, replacing the unknown biases by their estimates and setting the noise
to zero gives the nominal kinematics.

\begin{convention}[nominal propagation, AGENTS.md F1]
\begin{equation}
  \hat a = a_m - \hat b_a,\quad \hat\omega = \omega_m - \hat b_g;\qquad
  \dot{\hat p} = \hat v,\quad \dot{\hat v} = \hat R_{wb}\,\hat a + g_w,\quad
  \dot{\hat q} = \tfrac12\,\hat q\qprod[0,\ \hat\omega],\quad \dot{\hat b}_a = \dot{\hat b}_g = 0 .
  \label{eq:nominal-ct}
\end{equation}
The attitude update is always multiplicative, $\hat q \gets \hat q\qprod\Exp_q(\hat\omega\Dt)$, followed by
normalization. The integration scheme for $p$ and $v$ is chosen and justified in Milestone 4 and logged
in \texttt{PROGRESS.md}.
\end{convention}

The equations form a cascade: attitude depends on the gyroscope only; velocity depends on attitude and the
accelerometer; position depends on velocity. Errors cascade the same way, which is the origin of the
$t^2$ and $t^3$ growth laws of \cref{ch:intro}.

% ------------------------------------------------------------------------------
\section{Exact solution for constant inputs}

Suppose $\hat\omega$ and $\hat a$ are constant over $[t_k, t_k + \Dt]$, which is what a zero-order hold of one
IMU sample assumes. Then the kinematics can be solved in closed form.

\begin{proposition}[Constant-input propagation]
\label{prop:exact-prop}
With $\hat\omega$, $\hat a$ constant over the step and $\phi = \hat\omega\Dt$,
\begin{align}
  \hat R_{k+1} &= \hat R_k\,\Exp(\phi), \\
  \hat v_{k+1} &= \hat v_k + g_w\Dt + \hat R_k\,\Dt\,\Jl(\phi)\,\hat a, \\
  \hat p_{k+1} &= \hat p_k + \hat v_k\Dt + \tfrac12 g_w\Dt^2 + \hat R_k\,\Dt^2\,N(\phi)\,\hat a,
\end{align}
where $N(\phi) = \sum_{j\ge0}\sx{\phi}^j/(j+2)! = \tfrac12 I + \frac{\theta - \sin\theta}{\theta^3}\sx{\phi} +
\frac{\theta^2/2 + \cos\theta - 1}{\theta^4}\sx{\phi}^2$ and $\theta = \|\phi\|$.
\end{proposition}
\begin{proof}
The attitude is \cref{prop:const-rate}. During the step $\hat R(t_k + s) = \hat R_k\Exp(\hat\omega s)$, so
$\hat v_{k+1} - \hat v_k = g_w\Dt + \hat R_k\int_0^\Dt\Exp(\hat\omega s)\,\dd s\,\hat a$, and
$\int_0^\Dt\Exp(\hat\omega s)\,\dd s = \Dt\int_0^1\Exp(u\phi)\,\dd u = \Dt\,\Jl(\phi)$ by \cref{thm:jr}. Integrating once more
gives the position, with $\int_0^\Dt\!\int_0^s\Exp(\hat\omega u)\,\dd u\,\dd s = \Dt^2N(\phi)$ by term-wise integration of
the exponential series, and the closed form follows from $\sx{\phi}^3 = -\theta^2\sx{\phi}$.
\end{proof}

For small $\phi$, $\Jl(\phi) \approx I + \tfrac12\sx{\phi}$ and $N(\phi) \approx \tfrac12 I + \tfrac16\sx{\phi}$. The
correction terms account for the rotation of the specific-force vector during the step. At 200~Hz and
1~rad/s, $\|\phi\| = 0.005$, so they change the velocity increment by about 0.25\%: small, but systematic.

% ------------------------------------------------------------------------------
\section{Numerical integration schemes}

Real inputs are not constant over a step: the IMU provides samples at $t_k$ and $t_{k+1}$, and the true signal
varies between them. Three families of schemes are common. All of them update the attitude
multiplicatively.

\paragraph{Euler with zero-order hold.} Use the sample at the start of the step for the whole step:
\begin{equation}
  \hat R_{k+1} = \hat R_k\Exp(\hat\omega_k\Dt), \quad
  \hat v_{k+1} = \hat v_k + (\hat R_k\hat a_k + g_w)\Dt, \quad
  \hat p_{k+1} = \hat p_k + \hat v_k\Dt + \tfrac12(\hat R_k\hat a_k + g_w)\Dt^2 .
\end{equation}
Local error $O(\Dt^2)$, global error $O(\Dt)$. The dominant error is essentially a \emph{half-sample time
shift}: using the sample from the start of the interval for the whole interval delays the signal by
$\Dt/2$ on average.

\paragraph{Midpoint (trapezoidal).} Use the average of the two bounding samples:
\begin{equation}
  \hat R_{k+1} = \hat R_k\Exp\big(\tfrac12(\hat\omega_k + \hat\omega_{k+1})\Dt\big),\quad
  \bar a_w = \tfrac12\big(\hat R_k\hat a_k + \hat R_{k+1}\hat a_{k+1}\big) + g_w,
\end{equation}
\begin{equation*}
  \hat v_{k+1} = \hat v_k + \bar a_w\Dt, \qquad \hat p_{k+1} = \hat p_k + \hat v_k\Dt + \tfrac12\bar a_w\Dt^2 .
\end{equation*}
Local error $O(\Dt^3)$, global $O(\Dt^2)$, at the cost of having to wait for the next sample before
completing the step.

\paragraph{Runge--Kutta 4.} RK4 evaluates the dynamics at the start, twice at the middle, and at the end of
the step. Its global error is $O(\Dt^4)$ \emph{if} the inputs at mid-step are known exactly. With an IMU
they are not, and must be interpolated from the bounding samples. Linear interpolation is only
second-order accurate, which caps the whole scheme at $O(\Dt^2)$.

\begin{figure}[t]
  \centering
  \includegraphics[width=0.86\textwidth]{fig_integration_order.pdf}
  \caption{Position error after 10~s of noise-free integration on a synthetic trajectory (2~m circle at
  0.8~rad/s with $\pm23^\circ$ roll oscillations), as a function of the IMU period. The measured slopes
  match the predicted global orders. RK4 reaches fourth order only with exact mid-step inputs, which a
  real IMU does not provide; with interpolated inputs it is no better than the midpoint rule.}
  \label{fig:integration-order}
\end{figure}

\Cref{fig:integration-order} confirms the predicted orders. Halving $\Dt$ divides the error by 2 for Euler and
by 4 for the midpoint rule, which is exactly the convergence-order test of Milestone 4. The magnitudes
matter as much as the slopes. At 200~Hz on this fairly dynamic trajectory, the Euler scheme accumulates
about a meter of error in 10~s, far more than the sensor noise would cause (\cref{fig:cov-growth}), because the
half-sample delay of the attitude creates a tilt error of order $\tfrac12\omega\Dt$, which leaks gravity. The
midpoint rule is three orders of magnitude better. On gentler motion the difference shrinks, and frequent
visual updates bound integration error in any case, but the integration error should be well below the
noise-induced error, or it becomes an unmodeled error that the filter cannot account for.

\begin{keyidea}[choosing a scheme]
The integration error is a deterministic, unmodeled error. The scheme should make it small compared with
the stochastic error the filter does model. Whichever scheme Milestone 4 adopts, document its order and
verify it with the halving test on a trajectory with closed-form ground truth.
\end{keyidea}

\paragraph{Coning and sculling.} When the angular velocity vector itself rotates within a step (coning
motion) or rotation and acceleration oscillate in phase (sculling), single-sample schemes rectify
high-frequency motion into a spurious drift. High-end inertial systems with sample rates in the kHz range
use multi-sample correction algorithms~\cite{groves2013principles}. At 200~Hz on a flying robot whose
vibrations are mostly filtered out by the sensor, these effects are small compared with MEMS noise.

% ------------------------------------------------------------------------------
\section{Timing}

\subsection{Propagating to an arbitrary time}

Images arrive between IMU samples. To clone the state at an image time $t_c \in (t_k, t_{k+1})$ (AGENTS.md F5), the
filter must propagate exactly to $t_c$: interpolate the IMU reading at $t_c$ linearly between samples $k$ and
$k+1$, propagate from $t_k$ to $t_c$ with the interpolated input, clone, and then propagate from $t_c$ to $t_{k+1}$.
The covariance follows the same split, with $\Dt$ replaced by the two partial intervals.

\subsection{Event ordering and gaps}

Measurements must be processed in timestamp order, and the state must be propagated to a measurement's time
before the measurement is applied (\cref{ch:fullalgo}). Milestone 3 fixes the tie rule: an IMU sample precedes a
measurement with an equal timestamp, so the state is propagated to that time first. Gaps in the IMU stream
(intervals larger than about $1.5$ nominal periods) must not be assumed away. $\Dt$ always comes from the
timestamps, and a long gap simply produces a long step with a correspondingly large $Q_d$.

% ------------------------------------------------------------------------------
\section{Initialization}

\subsection{From ground truth}

The project initializes the nominal state from the first ground-truth sample at or after the start of the
stream (Milestone 4): $p$, $v$, $q$, $b_a$ and $b_g$ all come from the ground-truth columns, mapped with the
column order of AGENTS.md C7. This isolates propagation from the initialization problem. The use of ground
truth for initialization is logged when the full VIO runs (Milestone 10b).

\subsection{Static initialization}

A real system initializes from its own sensors. If the vehicle is at rest for about a second,
\begin{equation}
  \bar a_m = \mathrm{mean}(a_m) \approx R_{wb}\T(-g_w) + b_a, \qquad \bar\omega_m = \mathrm{mean}(\omega_m) \approx b_g .
\end{equation}
The gyro bias is the mean rate. Neglecting $b_a$, the unit vector $u_b = \bar a_m/\|\bar a_m\|$ is the world
``up'' direction expressed in the body frame, $u_b = R_{wb}\T e_z$. Any rotation with $R_{wb}u_b = e_z$ is consistent
with it: rotating about the axis $u_b\times e_z$ by the angle between the two vectors (computed with
$\atantwo$) gives one, and composing with any rotation about $e_z$ gives all the others. Roll and pitch are
determined; yaw is arbitrary, and is set to zero (or to the ground-truth yaw for evaluation).

\begin{pitfall}[tilt and accelerometer bias are indistinguishable at rest]
A horizontal accelerometer bias $b_a$ produces the same reading as a tilt of $\|b_a\|/g$: a bias of
$10~\mathrm{mm/s^2}$ looks exactly like a tilt of $1$~mrad. At rest the filter cannot separate them, and the initial
covariance must say so (large $\sigma_{b_a}$, correlated with the tilt). They become separable once the vehicle
rotates, because the bias rotates with the body while gravity does not.
\end{pitfall}

Initial covariances should reflect real uncertainty: a few milliradians in roll and pitch, a large or
unknown yaw, the turn-on bias range from the datasheet, and a small but non-zero velocity uncertainty at rest.
\emph{Dynamic} initialization from a moving start, as in VINS-Mono~\cite{qin2018vins}, first solves a
vision-only structure-from-motion problem and then aligns it with preintegrated IMU data to recover scale,
gravity and velocity.

% ------------------------------------------------------------------------------
\section{Sidebar: IMU preintegration}

An optimization-based estimator may relinearize a window of states many times. Re-integrating every IMU
sample between keyframes $i$ and $j$ each time the estimate of state $i$ changes would be wasteful.
\emph{Preintegration}~\cite{lupton2012visual, forster2017preintegration} integrates the measurements in the body frame
of keyframe $i$, so that the result does not depend on the state at $i$:
\begin{equation}
  \Delta R_{ij} = \prod_{k=i}^{j-1}\Exp(\hat\omega_k\Dt), \quad
  \Delta v_{ij} = \sum_{k=i}^{j-1}\Delta R_{ik}\,\hat a_k\Dt, \quad
  \Delta p_{ij} = \sum_{k=i}^{j-1}\Big(\Delta v_{ik}\Dt + \tfrac12\Delta R_{ik}\,\hat a_k\Dt^2\Big).
\end{equation}
These relate the two keyframe states through
$R_j = R_i\Delta R_{ij}$, $v_j = v_i + g\Delta t_{ij} + R_i\Delta v_{ij}$ and
$p_j = p_i + v_i\Delta t_{ij} + \tfrac12g\Delta t_{ij}^2 + R_i\Delta p_{ij}$. The preintegrated quantities depend on the
bias estimate, and first-order bias Jacobians avoid re-integration when the bias estimate changes. A filter
like the MSCKF does not need preintegration, because it integrates each sample exactly once.

\begin{projectlink}[Milestone 4]
Exit criteria: a stationary input keeps the state constant to $10^{-9}$ for 60~s; a constant rate reproduces
$q_0\qprod\Exp_q(\omega t)$; a circular trajectory converges at the predicted order; the quaternion norm stays
within $10^{-12}$ of one after $10^6$ steps; propagation does not allocate; and 10~s of dead reckoning on
MH\_01 produce no NaN.
\end{projectlink}

\begin{takeaways}
\begin{itemize}
  \item $\dot p = v$, $\dot v = R\hat a + g$, $\dot q = \tfrac12 q\qprod[0,\hat\omega]$; attitude updates are multiplicative.
  \item For constant inputs the solution is exact: $R\Exp(\phi)$, $v + g\Dt + R\Dt\Jl(\phi)\hat a$, and a similar
        position formula.
  \item Euler is first order (a half-sample delay), midpoint second order. RK4 is limited to second order by
        input interpolation.
  \item Integration error is unmodeled error; keep it well below sensor-noise error and verify the order.
  \item Propagate exactly to measurement times by splitting the IMU interval.
  \item At rest, gravity gives roll and pitch, the mean gyro gives $b_g$, and tilt and horizontal $b_a$ are
        indistinguishable.
\end{itemize}
\end{takeaways}

\begin{exercises}
  \item Verify \cref{prop:exact-prop} numerically by comparing with an ODE solver at tight tolerance.
  \item Show that the Euler scheme applied to a constant rotation rate is exact for the attitude, but not for
        the velocity, and quantify the velocity error per step in terms of $\|\hat\omega\|\Dt$.
  \item Reproduce \cref{fig:integration-order}. Then add gyro noise and plot the noise-induced position error
        on the same axes. At what IMU rate does integration error fall below noise error for each scheme?
  \item Implement the static initializer and test it on synthetic data with a known tilt. Add a horizontal
        accelerometer bias and show that it biases the estimated tilt by $\|b_a\|/g$.
  \item Why must the propagation to an image timestamp interpolate the IMU sample instead of using the
        nearest sample? Estimate the attitude error from using the nearest sample at 1~rad/s and 200~Hz.
\end{exercises}
